单位圆 \(S^1=\{(x,y):x^2+y^2=1\}\) 是一个一维流形。去掉北极点 \((0,1)\) 后可用一个图表 \(\varphi_1(x,y)=\arctan(y/x)\) 参数化；去掉南极点 \((0,-1)\) 后用 \(\varphi_2(x,y)=\arctan(y/x)\)（另一分支）覆盖剩余部分。两个图表的交是去掉南北极的圆，坐标转换是平移——这就是``局部像 \(\R\)，全局粘合''。

Part 13 的降维单元反复说``数据躺在低维流形上''，Part 15 的动力系统说``稳定流形与不稳定流形''。本单元把这个词变成可检验的数学对象：先理解流形是什么（局部像欧氏空间），再理解怎样在流形上移动（切空间），然后理解怎样测量距离（黎曼度量），最后理解怎样研究流形上的函数（光滑映射）。

读完应能判断一段关于流形的陈述是定义、定理还是建模假设，也能在 Part 13 的降维与 Part 15 的动力系统中，把``流形''``切空间''``测地距离''从比喻变成可计算的句子。

\subsection{一个具体算例}

\(S^1\) 在点 \(p=(\cos\theta_0,\sin\theta_0)\) 处的切空间是过 \(p\) 与圆心连线垂直的直线，基向量为 \((-\sin\theta_0,\cos\theta_0)\)。若在 \(S^1\) 上用诱导黎曼度量 \(ds^2=d\theta^2\)（取 \(R=1\)），弧长即角度差。测地线是大圆弧：从 \(\theta=0\) 到 \(\theta=\pi/3\) 的测地距离为 \(\pi/3\)，而弦长只有 \(2\sin(\pi/6)=1\)——在弯曲空间里``直线''不一定给出最短的弦。

\subsection{怎么读这一章}

先把流形理解成``局部能用欧氏坐标描述的弯曲空间''。坐标图只负责定位，切空间恢复可移动方向，黎曼度量才赋予长度与角度，最后再看函数和映射怎样在这些局部坐标间保持一致。

\subsection{本章篇目}

以下篇目按概念依赖排列；若从具体问题进入，也可以先打开最接近当前困惑的一篇，再沿篇内链接回补。

\begin{itemize}
\tightlist
\item \hyperref[sec:12-Real-Analysis-Support/08-Manifold-Basics/01]{``流形是局部像欧氏空间的集合''}；
\item \hyperref[sec:12-Real-Analysis-Support/08-Manifold-Basics/02]{``切空间把``方向''还给弯曲空间''}；
\item \hyperref[sec:12-Real-Analysis-Support/08-Manifold-Basics/03]{``黎曼度量让弯曲空间有长度''}；
\item \hyperref[sec:12-Real-Analysis-Support/08-Manifold-Basics/04]{``流形上的函数与映射''}。
\end{itemize}
